\section*{Chapter \ref{ch:b1:nernst} --- Redox Equilibria and the Nernst Equation}

\begin{solution}{exo:b1:nernst:1}
N in \ce{NH3}: $-$III; in \ce{NO3-}: $+$V. S in \ce{SO4^2-}: $+$VI. O in
\ce{H2O2}: $-$I. Cr in \ce{Cr2O7^2-}: $+$VI. Cl in \ce{ClO-}: $+$I. Fe in
\ce{Fe3O4}: $+8/3$ on average (one iron(II) and two iron(III)). C in
\ce{CH3OH}: $x + 4(+1) + (-2) = 0$, $-$II.
\end{solution}

\begin{solution}{exo:b1:nernst:2}
\ce{Cr2O7^2- + 14H+ + 6e- -> 2Cr^3+ + 7H2O};
\ce{H2O2 + 2H+ + 2e- -> 2H2O};
\ce{O2 + 2H+ + 2e- -> H2O2};
\ce{SO4^2- + 4H+ + 2e- -> SO2 + 2H2O}.
\end{solution}

\begin{solution}{exo:b1:nernst:3}
$E = -0.76 + 0.030\log[\ce{Zn^2+}]$; $E = 1.36 + 0.030\log(p(\ce{Cl2})/[\ce{Cl-}]^2)$;
$E = 1.23 + 0.015\log(p(\ce{O2})[\ce{H+}]^4)$; $E = 0.27 - 0.059\log[\ce{Cl-}]$
(pressures in bar). In \qty{0.010}{mol/L} zinc sulfate:
$E = -0.76 + 0.030 \times (-2) = -0.82$~V.
\end{solution}

\begin{solution}{exo:b1:nernst:4}
$E(\ce{Zn}) = -0.76 + 0.030\log 0.10 = -0.79$~V, $E(\ce{Cu}) = 0.34 +
0.030\log 0.010 = 0.28$~V. Voltage $0.28 - (-0.79) = 1.07$~V. Zinc, at the
lower potential, is the negative pole, where oxidation takes place: the
\omterm{def:b1:nernst:cell}{anode}.
\end{solution}

\begin{solution}{exo:b1:nernst:5}
\ce{MnO4- + 5Fe^2+ + 8H+ -> Mn^2+ + 5Fe^3+ + 4H2O}, $n = 5$,
$\log K = 5(1.51 - 0.77)/0.059 = 63$: \omterm{def:b1:extent-q-and-k:quantitative}{quantitative} and fast, a classic
titration (permanganate is its own indicator).
\end{solution}

\begin{solution}{exo:b1:nernst:6}
Metals whose $E^\circ$ is below 0 react: Mg, Zn, Fe, Ni, Pb (lead only
slowly: the gap is small); Cu and Ag do not. \ce{Mg + 2H+ -> Mg^2+ + H2}, $\log K = 2 \times 2.36/0.059 = 80$;
\ce{Fe + 2H+ -> Fe^2+ + H2}, $\log K = 2 \times 0.41/0.059 = 13.9$.
\end{solution}

\begin{solution}{exo:b1:nernst:7}
$n = 2$, $\log K = 2(0.53 - 0.02)/0.059 = 17$: \omterm{def:b1:extent-q-and-k:quantitative}{quantitative}.
\end{solution}

\begin{solution}{exo:b1:nernst:8}
$3E^\circ(\ce{Fe^3+}/\ce{Fe}) = 2(-0.41) + 0.77$, $E^\circ = -0.02$~V.
\ce{2Fe^3+ + Fe -> 3Fe^2+}: \omterm{def:b1:nernst:couple}{oxidant} \ce{Fe^3+} (0.77) above \omterm{def:b1:nernst:couple}{reductant}
\ce{Fe} ($-0.41$), $n = 2$, $\log K = 2(0.77 + 0.41)/0.059 = 40$. Iron
filings keep a solution of iron(II) from turning into iron(III).
\end{solution}

\begin{solution}{exo:b1:nernst:9}
\ce{O2 + 4H+ + 4e- -> 2H2O}: $E = 1.23 + \frac{0.059}{4}\log [\ce{H+}]^4 =
1.23 - 0.059\,\mathrm{pH}$. \ce{2H+ + 2e- -> H2}: $E = \frac{0.059}{2}\log
[\ce{H+}]^2 = -0.059\,\mathrm{pH}$. The difference is $1.23$~V at any pH:
the voltage of the fuel cell does not depend on the pH.
\end{solution}

\begin{solution}{exo:b1:nernst:10}
$E^\circ(\ce{Cu+}/\ce{Cu}) = 0.52 > E^\circ(\ce{Cu^2+}/\ce{Cu+}) = 0.16$: the
\omterm{def:b1:nernst:couple}{oxidant} \ce{Cu+} of the upper couple reacts with the \omterm{def:b1:nernst:couple}{reductant} \ce{Cu+} of
the lower. $\log K = (0.52 - 0.16)/0.059 = 6.1$, $K =
[\ce{Cu^2+}]/[\ce{Cu+}]^2$, so $[\ce{Cu+}] = \sqrt{0.10/10^{6.1}} =
\qty{2.8e-4}{mol/L}$.
\end{solution}

\begin{solution}{exo:b1:nernst:11}
$U = 0.059\log(0.10/\num{1.0e-3}) = 0.118$~V. The positive pole is the
silver in the concentrated solution, where \ce{Ag+} is reduced; in the
dilute solution silver is oxidised. The concentrated solution becomes
more dilute and the dilute one more concentrated; the voltage falls and
reaches zero when the two concentrations are equal.
\end{solution}

\begin{solution}{exo:b1:nernst:12}
\ce{H2O2} is the \omterm{def:b1:nernst:couple}{oxidant} of \ce{H2O2}/\ce{H2O} (1.76~V) and the \omterm{def:b1:nernst:couple}{reductant}
of \ce{O2}/\ce{H2O2} (0.69~V): it reacts with itself, \ce{2H2O2 -> 2H2O + O2},
$\log K = 2(1.76 - 0.69)/0.059 = 36$. The reaction is favoured but very
slow without a \omterm{def:b1:elementary-steps:catalyst}{catalyst} (\cref{ch:b1:elementary-steps}); clean bottles,
cold and dark storage keep it slow.
\end{solution}

\begin{solution}{pb:b1:nernst:1}
\textbf{1.} Ag: 0, $+$I, $+$I; Cl: $-$I in both.
\textbf{2.} \ce{AgCl(s) + e- -> Ag(s) + Cl-}.
\textbf{3.} $E = E^\circ(\ce{AgCl}/\ce{Ag}) - 0.059\log[\ce{Cl-}]$.
\textbf{4.} The two solids have \omterm{def:b1:extent-q-and-k:activity}{activity} 1; only the chloride ion is in
solution.
\textbf{5.} $E^\circ = 0.22$~V.
\textbf{6.} $0.22 + 0.118 = 0.34$~V.
\textbf{7.} \ce{Pt} $|$ \ce{H2} (\qty{1}{bar}) $|$ \ce{H+} (\omterm{def:b1:extent-q-and-k:activity}{activity} 1)
$\|$ \ce{Cl-} (\qty{0.10}{mol/L}) $|$ \ce{AgCl} $|$ \ce{Ag}.
\textbf{8.} The silver electrode: $0.22 + 0.059 = 0.28$~V $> 0$.
\textbf{9.} $0.28$~V.
\textbf{10.} \ce{2AgCl + H2 -> 2Ag + 2H+ + 2Cl-}.
\textbf{11.} $\log K = 2 \times 0.22/0.059 = 7.5$.
\textbf{12.} From the platinum (negative pole, where \ce{H2} is oxidised)
to the silver.
\textbf{13.} $0.27 - 0.059\log 1.0 = 0.27$~V.
\textbf{14.} $E(\ce{Ag}) = 0.22 + 0.177 = 0.40$~V $> 0.27$: the silver
electrode is positive, $U = 0.13$~V.
\textbf{15.} $U = 0.22 - 0.059\log[\ce{Cl-}] - 0.27 = -0.05 -
0.059\log[\ce{Cl-}]$ (volts).
\textbf{16.} $\log[\ce{Cl-}] = -(0.115 + 0.05)/0.059 = -2.80$,
$[\ce{Cl-}] = \qty{1.6e-3}{mol/L}$, that is \qty{57}{mg/L}.
\textbf{17.} $\Delta c/c = \ln 10 \times 0.001/0.059 = 0.039$: about
\qty{4}{\percent} per millivolt.
\textbf{18.} When the chloride of the sample becomes comparable to what
silver chloride itself releases, $\sqrt{K_s} = \qty{1.3e-5}{mol/L}$;
the electrode is used well above, from about \qty{1e-4}{mol/L}.
\textbf{19.} $E = 0.80 + 0.059\log[\ce{Ag+}]$.
\textbf{20.} $[\ce{Ag+}] = K_s/[\ce{Cl-}]$.
\textbf{21.} $E = 0.80 + 0.059\log K_s - 0.059\log[\ce{Cl-}] = 0.80 -
0.059\,\mathrm{p}K_s - 0.059\log[\ce{Cl-}]$.
\textbf{22.} $E^\circ(\ce{AgCl}/\ce{Ag}) = E^\circ(\ce{Ag+}/\ce{Ag}) -
0.059\,\mathrm{p}K_s$ (\cref{prop:b1:nernst:second-kind}).
\textbf{23.} $0.80 - 0.059 \times 9.75 = 0.225$~V, in agreement with the
table (0.22~V), computed independently from Gibbs energies.
\textbf{24.} $E^\circ(\ce{AgCl}/\ce{Ag}) = $ \textbf{0.22~V}.
\end{solution}
